\documentclass[10pt,a4paper]{amsart}
\usepackage[margin=19mm]{geometry}
\usepackage{amsmath,amssymb,mathtools}
\usepackage[scr=rsfs,cal=euler]{mathalfa}
\usepackage{xfrac}
\usepackage[unicode,hidelinks,bookmarks=false]{hyperref}
\newcommand{\ie}{i.e.\ }
\newcommand{\cf}{cf.\ }
\newcommand{\resp}{respectively\ }
\newcommand{\Subsection}{\subsection}
\providecommand{\nobreakdash}{\nobreak\hbox{-}}
\setlength{\emergencystretch}{2em}
\multlinegap0pt
\usepackage{amssymb}
\usepackage{mathtools}
\usepackage[shortlabels]{enumitem}
\newtheorem{lemma}{Lemma}
\newcommand{\R}{\mathbb R}
\newcommand{\cL}{\mathcal L}
\newcommand{\dual}[2]{\langle #1,#2\rangle}
\title{H\"older Stability from Exact Uniqueness for Finite-Dimensional Analytic Inverse Problems}
\author{C\u{a}t\u{a}lin I. C\^arstea}
\date{Source: arXiv:2605.06354v1 (CC BY 4.0)}
\begin{document}
\maketitle

\noindent\emph{Source and licence.} H\"older Stability from Exact Uniqueness for Finite-Dimensional Analytic Inverse Problems. Version arXiv:2605.06354v1 (CC BY 4.0), distributed by arXiv under the Creative Commons Attribution 4.0 International licence, http://creativecommons.org/licenses/by/4.0/, as recorded in the arXiv licence metadata for this exact version and shown on the abstract page https://arxiv.org/abs/2605.06354v1. Full licence notice: \texttt{licenses/arxiv-2605.06354-CC-BY-4.0.txt}.\par\smallskip

\noindent\textit{Editorial scope.} Counted unit: the article's lemma lem:elastic-analytic (source0.tex lines 1141--1147). The article's complete original proof of it is delivered with it (source0.tex lines 1149--1243, one \texttt{proof} environment of 95 lines). No mathematics is written, repaired or re-derived: the statement and the proof are the article's own lines, in the article's own notation, and no retained line is substituted or re-flowed.  Retained uncounted as the source-local context the proof needs: lines 976--976.  Nothing was omitted from any retained span, so the package carries no omission record.  No retained line contains a citation, so the wrapper carries no bibliography and no bibliography entry.  No retained line contains a figure environment, an image-inclusion command or drawing code, so no image is delivered and the package declares no asset dependency.  Editorial formatting only, in addition: an AMS \texttt{amsart} preamble that restates the article's own declarations and macros used by the retained lines, namely the environments \texttt{lemma} on separate counters, together with the standard packages those lines need.  The retained environments print in this document as Lemma 1 in order of appearance, whereas the article prints the same items with the numbering of its own full document.  The complete native source of the article as distributed in the arXiv e-print for this exact version is bundled unchanged as \texttt{sources/200012/source0.tex}.
\subsection{Piecewise homogeneous elasticity and localized D-N data}\label{subsec:elasticity-data}

\begin{lemma}[Analyticity of the localized elasticity D-N map]\label{lem:elastic-analytic}
The map
\[
        \mathbb C\mapsto\Lambda_{\mathbb C}^\Gamma
\]
is real analytic from $U_{\rm el}$ into $\cL(X,X^*)$.
\end{lemma}

\begin{proof}
We fix one extension operator and solve only for the zero-boundary correction.

Let
\[
        V_0=H^1_0(\Omega;\R^3).
\]
Choose a bounded right inverse of the trace map
\[
        E:H^{1/2}(\partial\Omega;\R^3)\to H^1(\Omega;\R^3),
\]
and, for $f\in X$, apply it to the zero extension of $f$ from $\Gamma$ to $\partial\Omega$.  We still denote the resulting bounded map $X\to H^1(\Omega;\R^3)$ by $E$.

For the symmetric-gradient form
\[
        a_{\mathbb C}(u,v)=
        \int_\Omega \mathbb C(x)\varepsilon(u):\varepsilon(v)\,dx,
\]
define
\[
        L_{\mathbb C}:V_0\to V_0^*,
        \qquad
        \dual{L_{\mathbb C}w}{v}=a_{\mathbb C}(w,v).
\]
The full-gradient formulation used in the statement of the inverse problem agrees with this symmetric-gradient form by the index calculation in Subsection \ref{subsec:elasticity-data}.  The map $\mathbb C\mapsto L_{\mathbb C}$ is linear: its matrix coefficients are finite linear combinations of fixed operators of the form
\[
        (w,v)\mapsto
        \int_{D_\alpha}\varepsilon_{ij}(w)\varepsilon_{k\ell}(v)\,dx .
\]
Strong convexity and Korn's inequality give coercivity on $V_0$, hence $L_{\mathbb C}$ is a bounded isomorphism for every $\mathbb C\in U_{\rm el}$.  Therefore
\[
        \mathbb C\mapsto L_{\mathbb C}^{-1}
\]
is real analytic into $\cL(V_0^*,V_0)$.

Now define two auxiliary operators.  First,
\[
        P_{\mathbb C}:X\to V_0^*,
        \qquad
        \dual{P_{\mathbb C}f}{v}=a_{\mathbb C}(Ef,v),
        \quad v\in V_0 .
\]
Second,
\[
        Q_{\mathbb C}:X\to X^*,
        \qquad
        \dual{Q_{\mathbb C}f}{g}=a_{\mathbb C}(Ef,Eg).
\]
Both $P_{\mathbb C}$ and $Q_{\mathbb C}$ depend linearly, hence analytically, on $\mathbb C$.  If $f\in X$, the solution with boundary value $f$ has the form
\[
        u_{\mathbb C}^f=Ef+w_{\mathbb C}^f,
        \qquad w_{\mathbb C}^f\in V_0.
\]
The weak equation $a_{\mathbb C}(u_{\mathbb C}^f,v)=0$ for every $v\in V_0$ gives
\[
        L_{\mathbb C}w_{\mathbb C}^f=-P_{\mathbb C}f,
        \qquad
        w_{\mathbb C}^f=-L_{\mathbb C}^{-1}P_{\mathbb C}f.
\]
Thus the solution operator is analytic after the fixed lift $Ef$ is separated off.

It remains to write the boundary traction pairing in terms of these operators.  For $z\in V_0$, define
\[
        P_{\mathbb C}^{\sharp}z\in X^*,
        \qquad
        \dual{P_{\mathbb C}^{\sharp}z}{g}
        =\dual{P_{\mathbb C}g}{z}
        =a_{\mathbb C}(Eg,z).
\]
The map $\mathbb C\mapsto P_{\mathbb C}^{\sharp}$ is again linear.  Since $u_{\mathbb C}^f$ solves the homogeneous equation, the D-N pairing may be computed with the extension $Eg$ of the boundary test $g$:
\[
\begin{aligned}
        \dual{\Lambda_{\mathbb C}^\Gamma f}{g}
        &=a_{\mathbb C}(u_{\mathbb C}^f,Eg) \\
        &=a_{\mathbb C}(Ef,Eg)
          -a_{\mathbb C}(L_{\mathbb C}^{-1}P_{\mathbb C}f,Eg) \\
        &=\dual{Q_{\mathbb C}f}{g}
          -\dual{P_{\mathbb C}^{\sharp}L_{\mathbb C}^{-1}P_{\mathbb C}f}{g}.
\end{aligned}
\]
Here $L_{\mathbb C}^{-1}P_{\mathbb C}f\in V_0$, so $P_{\mathbb C}^{\sharp}$ applies to it.  The last equality uses the symmetry of $a_{\mathbb C}$ to rewrite
\[
        a_{\mathbb C}(L_{\mathbb C}^{-1}P_{\mathbb C}f,Eg)
        =a_{\mathbb C}(Eg,L_{\mathbb C}^{-1}P_{\mathbb C}f).
\]
Therefore, as an operator $X\to X^*$,
\[
        \Lambda_{\mathbb C}^\Gamma
        =Q_{\mathbb C}
        -P_{\mathbb C}^{\sharp}L_{\mathbb C}^{-1}P_{\mathbb C}.
\]
This formula expresses the D-N map using only linear functions of $\mathbb C$, the analytic inverse $L_{\mathbb C}^{-1}$, and bounded bilinear composition of operators.  Hence $\mathbb C\mapsto\Lambda_{\mathbb C}^\Gamma$ is real analytic in $\cL(X,X^*)$.

The formula does not depend on the auxiliary choice of $E$.  If the test extension $Eg$ is changed by an element of $V_0$, the extra term pairs to zero with $u_{\mathbb C}^f$ by the weak equation.  If the lift of $f$ is changed, the zero-boundary correction changes by the opposite zero-boundary function, so the resulting solution $u_{\mathbb C}^f$ is unchanged by uniqueness.
\end{proof}

\end{document}
